\begin{comment}
=== AI-INSTRUCTIONS (NEVER DELETE OR ALTER THIS BLOCK) ===
1. VERBATIM RETENTION: Always include this exact comment block at the top of your output.
2. LHS STYLE: Explanations on the Left-Hand Side must be colloquial, punchy, and use "footnotesize".
3. DRY/KISS: Always tie the code back to SE principles (like DRY, KISS, YAGNI, SRP) where applicable.
4. SPATIAL LIMITS: Keep the minted code to ~40 lines maximum.
5. REVIEW NOTES:
   - [Tim's Note]: Make sure to emphasize why we aren't using Pandas (YAGNI).
      - [Tim's Note]: Keep the magic characters (+, -, !) explanation prominent.

      - [T note]: use of clone.
         - [T note]:  mention this is active learninger (best results when learners pick their own data.
                         by focusing on most informative examples, we can dodge noise and irrelvancies and superflous signnals and learn from
                                         vastly less data.
   - [T note]:  there are many constructions of first code sample that would confuse/mislead newbies. e.g.
               comprehensions. comprejensions with 
               conditionals.
               type
               all such idoms need to be teased out.
               maybe takes more than 1 page


      ==========================================================
\end{comment}

% ==========================================
% SLIDE 1: Core Architecture (Code)
% ==========================================
\begin{frame}[fragile]{The EZR Loop: Herding the System}
\begin{columns}[T]
  \begin{column}{0.4\textwidth}
    \footnotesize 
    \textbf{Lines 314-345:}\\
    This is an \textbf{Active Learner}. We get the absolute best results when learners pick their own data! By focusing on the most informative examples, we dodge noise, ignore superfluous signals, and learn from vastly less data.

    \vspace{0.2cm}
    \rc{1} \textbf{Bootstrap:} Tiny random sample. Software spaces are mostly deserts anyway.
    
    \vspace{0.15cm}
    \rc{2} \textbf{The Clone Wars:} Divide seen data into \texttt{best} vs \texttt{rest} by heavily reusing our \texttt{clone} function (\textbf{DRY}).
    
    \vspace{0.15cm}
    \rc{3} \textbf{Cheap Guessing:} Use fast geometry to find points close to \texttt{best}, far from \texttt{rest}. No expensive oracle calls! (\textbf{KISS})
    
    \vspace{0.15cm}
    \rc{4} \textbf{Secret Sauce:} Cap \texttt{best} size at $\sqrt{N}$. Dump overflows to \texttt{rest}.
  \end{column}
  
  \begin{column}{0.6\textwidth}
\begin{minted}{python}
# ---- 7. Active learning ----
def acquire(d, score=acquireWithBayes, label=lambda _,r:r):
  """Pick what unlabelled row to evaluate next."""
  rows = d.rows[:]
  shuffle(rows)
  
  # 1. Initialize labelled vs unlabelled
  lab   = clone(d, rows[:the.learn.start]) |\rc{1}|
  unlab = rows[the.learn.start:][:the.few]
  lab.rows.sort(key=lambda r: disty(lab, label(d,r)))
  
  # 2. Split seen data
  n = sqrt(len(lab.rows))
  best = clone(d,lab.rows[:int(n)]) |\rc{2}|
  rest = clone(d,lab.rows[int(n):])
  
  cursor = 0
  for _ in range(the.learn.budget):
    for j in range(len(unlab)):  
      idx = (cursor + j) % len(unlab)  
      if score(lab, best, rest, unlab[idx]) < 0: |\rc{3}|
        # 3. Label and add
        add(lab, add(best, label(d, unlab.pop(idx))))
        
        # 4. Prune the 'best' cluster
        if len(best.rows) > sqrt(len(lab.rows)): |\rc{4}|
          best.rows.sort(key=lambda r: disty(lab,r))
          add(rest, sub(best, best.rows[-1]))
          
        cursor = idx  
        break
  return lab
\end{minted}
  \end{column}
\end{columns}
\end{frame}

% ==========================================
% SLIDE 2: Scripting Notes
% ==========================================
\begin{frame}[fragile]{Scripting Notes: Teasing Out Python Idioms}
  \vspace{0.2cm}
  Python has idioms that can easily trip up newbies. Let's break down the tricks used in the \texttt{acquire} function.

  \vspace{0.3cm}
  \textbf{Idiom 1: Shallow Copying with Slices}
  \begin{minted}[fontsize=\footnotesize]{python}
  rows = d.rows[:]
  \end{minted}
  We want to shuffle our rows to randomize them, but \texttt{shuffle} operates in-place. If we didn't add \texttt{[:]}, we would destroy the original ordering in our main \texttt{Data} object! The slice guarantees a fast, safe copy.
  
  \vspace{0.3cm}
  \textbf{Idiom 2: Anonymous Default Functions (Lambdas)}
  \begin{minted}[fontsize=\footnotesize]{python}
  def acquire(d, score=acquireWithBayes, label=lambda _,r:r):
  \end{minted}
  We pass a \texttt{lambda} as a default argument. The underscore \texttt{\_} tells the reader "we don't care about the first argument here." By default, it just returns the row \texttt{r} unchanged.
\end{frame}

% ==========================================
% SLIDE 3: Scripting Notes & SE Principles
% ==========================================
\begin{frame}[fragile]{Scripting Notes \& SE Principles}
  
  \vspace{0.2cm}
  \textbf{Idiom 3: Popping By Index}
  \begin{minted}[fontsize=\footnotesize]{python}
  add(lab, add(best, label(d, unlab.pop(idx))))
  \end{minted}
  \texttt{pop(idx)} does two things at exactly the same time: it returns the value at that index to be evaluated, and it removes it from the \texttt{unlab} pool so we never evaluate it again. 

  \vspace{0.4cm}
  \begin{block}{\textbf{DRY} (Don't Repeat Yourself)}
  Look at how we construct the \texttt{lab}, \texttt{best}, and \texttt{rest} pools. We don't write custom cluster-management logic. We just use \texttt{clone(d, subset)}. It builds an entirely new \texttt{Data} object with identical schema but new rows. 
  \end{block}

  \begin{block}{\textbf{KISS} (Keep It Simple, Stupid)}
  Calling an external software system (the "Oracle") is computationally expensive. Running a quick geometric \texttt{score} calculation in memory is virtually free. We keep it simple by guessing with math before we verify with the system.
  \end{block}
\end{frame}
